\section{Related Work}
\label{sec:related}

\subsection{Spurious Correlations and Group Robustness}

The study of spurious correlations in deep learning traces to observations that models exploit dataset biases rather than learning causal features \citep{Geirhos2020Shortcut, Zhang2017Rethinking}. \citet{Sagawa2020GroupDRO} formalize this via group robustness: in Waterbirds and CelebA, a model achieving high average accuracy still fails badly on minority groups whose spurious attribute contradicts the majority correlation. Their Group DRO method explicitly minimizes worst-group loss, requiring group annotations at training time. Our work is complementary: we characterize how much spurious content enters frozen representations \textit{before} any downstream training, as a function of pretraining paradigm --- without assuming group labels are available.

\citet{Kirichenko2022DFR} demonstrate that ERM features are sufficient for robustness: only the classifier head needs retraining on a small group-balanced set (Deep Feature Reweighting, DFR), achieving near-oracle performance on Waterbirds. This result motivates studying frozen representation quality directly, as we do. However, DFR focuses entirely on supervised ERM and does not compare to SSL paradigms --- leaving open whether ERM's spurious feature content is typical or atypical among pretraining objectives.

\citet{Izmailov2022Feature} provide a broader study of feature learning under spurious correlations, including comparisons with DINO on Waterbirds. Their analysis is closest to ours in scope, but does not use a controlled spurious/task ratio metric and does not perform a 4-paradigm comparison with Bonferroni-corrected pairwise tests. We provide this systematic comparison.

\subsection{Self-Supervised and Contrastive Representation Learning}

Contrastive SSL methods \citep{Chen2020SimCLR, He2020MoCo} learn representations by maximizing agreement between two augmented views of the same image. DINO \citep{Caron2021DINO} uses self-distillation with a momentum teacher, producing emergent class-discriminative attention maps. BarlowTwins \citep{Zbontar2021BarlowTwins} optimizes redundancy-reduction across cross-correlation matrices. All four paradigms use ResNet-50 in their standard evaluation protocols, making controlled comparison feasible.

These methods are primarily evaluated on ImageNet linear probe accuracy \citep{Chen2020SimCLR, Caron2021DINO}, which measures task-relevant encoding but says nothing about spurious feature content. The key gap is that \textit{no prior work measures spurious attribute probe accuracy as a function of pretraining paradigm using a controlled backbone and group-annotated datasets.}

\citet{Robinson2021Contrastive} show that contrastive learning can encode shortcut features, particularly when spurious attributes are stable across augmentations. \citet{Wen2021Feature} analyze what features contrastive objectives encode as a function of augmentation design. These works establish that SSL is not spurious-feature-free --- but they do not compare to supervised baselines under controlled conditions. We do, and find that supervised ERM encodes spurious features \textit{more} strongly, not less.

\subsection{Distribution Shift Benchmarks}

WILDS \citep{Koh2021WILDS} provides ten real-world distribution shift datasets with standardized splits, including Waterbirds and CelebA under the spurious correlation framing. DomainBed \citep{Gulrajani2021DomainBed} benchmarks 20+ domain generalization algorithms on seven datasets, showing that careful ERM tuning matches or surpasses specialized methods. Both evaluate \textit{fine-tuned} models, making it impossible to attribute performance differences to the pretraining representation versus the fine-tuning procedure. We isolate the pretraining effect by freezing the backbone throughout.

\subsection{Augmentation and Feature Learning}

\citet{Shah2020Simplicity} characterize SGD's ``simplicity bias'' --- the tendency to prefer low-complexity predictors even when more complex task-relevant features exist. This connects to our finding that ERM encodes spurious features aggressively: spurious attributes (background texture) are often simpler and more predictive than task-relevant features (bird morphology) in biased training distributions. \citet{Zimmermann2021Contrastive} show theoretically that contrastive learning approximately inverts the data generating process, recovering latent structure. Our empirical finding --- that label-agnostic contrastive SSL encodes spurious features less than label-supervised ERM --- is consistent with this theoretical framework when spurious attributes are not in the true latent structure.

\textbf{Our position:} We synthesize and extend these threads by providing, to the best of our knowledge, the first controlled 4-paradigm $\times$ 2-dataset comparison of spurious attribute linear probe accuracy on frozen ResNet-50 representations, using a spurious/task ratio metric that decouples paradigm effects from absolute accuracy differences across datasets.
